\section{Incidence spaces, parameters and retained distances}
For an $s$-subset $\Lambda\subseteq\Omega$, let $w_\Lambda$ be the
indicator of the $k$-subsets containing $\Lambda$; put
$W_s=\langle w_\Lambda:|\Lambda|=s\rangle$ and
$E(W_s)=\langle w_\Lambda-w_{\Lambda'}\rangle$.
These are the spaces of source equation~(4). The alternating transversal
word associated with $k$ disjoint pairs, from source equation~(5),
belongs to $W_s^\perp$ for $s<k$: for any fixed $s$-set, either no
transversal contains it, or at least one pair is unspecified and its
two choices contribute opposite signs. Thus $W_\Pi\subseteq W_s^\perp$
over every characteristic. This elementary cancellation is all that is
needed of that notation here.

\begin{proposition}\label{prop:point-incidence}
Suppose $1\leq k<n$. Then $\dim W_1=n$ if $p\nmid k$, and $n-1$
if $p\mid k$. In the binary case set
\[
 b_r=\sum_{i\text{ odd}}\binom ri\binom{n-r}{k-i},\qquad
 \epsilon=[k\text{ is even}].
\]
The full Hamming weight enumerator, including repeated weights and the zero
word, is
\begin{equation}\label{eq:W1-enumerator}
 \sum_{c\in W_1}z^{\wt(c)}
     =2^{-\epsilon}\sum_{r=0}^{n}\binom nr z^{b_r}.
\end{equation}
\end{proposition}
\begin{proof}
Map $(a_x)_{x\in\Omega}$ to $\sum_xa_xw_x$.
If this is zero, compare its coordinates on $T\cup\{x\}$ and
$T\cup\{y\}$, where $|T|=k-1$ and $x,y\notin T$. Such a $T$ exists
for every $x\ne y$ because $k<n$. It follows that $a_x=a_y$.
Hence the kernel consists exactly of constant coefficient vectors
$a\mathbf j$ with $ka=0$, proving the dimension.
In characteristic two a coefficient support $S$ of size $r$ gives a
nonzero coordinate at $X$ precisely when $|S\cap X|$ is odd, and the
count is $b_r$. Each word has exactly $2^\epsilon$ coefficient
preimages. Dividing the sum over all coefficient supports by that
constant fiber size proves \eqref{eq:W1-enumerator}.
\end{proof}
This recovers and qualifies source Result~4. The dimension statement
requires $k<n$, as in all its graph applications. Its displayed weight
values follow from $b_{n-r}=b_r$ for even $k$ and
$b_{n-r}=\binom nk-b_r$ for odd $k$.
Its unqualified claim of $\binom nr$ words of weight $b_r$ must instead
respect fibers and collisions. For example, when $(n,k)=(8,2)$,
weight $16$ has $\binom84/2=35$ words, not $70$.
Formula~\eqref{eq:W1-enumerator} also includes the constant word when
$k$ is odd. These corrections refer to the statement printed in F017;
they are not assertions about its cited predecessor.

\subsection{The complete quoted parameter list}
Source Result~3 quotes earlier code classifications, principally
Key--Moori--Rodrigues and Fish--Key--Mwambene, listed as source references
[11,14,16,17,21--23]. The following table preserves every listed case.
For $k=3$ take $n\geq7$ and $N=\binom n3$; it gives the reflexive space,
its dimension and distance, and those of its adjacency dual.
\begin{center}\small
\begin{tabular}{lllll}\toprule
condition&$RC_r=C_r^\perp$&$\dim RC_r$&$d(RC_r)$&$d(C_r)$\\\midrule
$r=0,p\geq5,n\equiv5\pmod p$&$W_2$&$\binom n2$&$n-2$&$\leq8$\\
$r=0,p=2,n\equiv1\pmod4$&$W_1+W_2$&$\binom n2$&$n-2$&$8$\\
$r=0,p=3,n\equiv5\pmod9$&$W_2+\langle\mathbf j\rangle$&$\binom n2$&$n-2$&$8$\\
$r=1,p=2,n\equiv0\pmod4$&$W_1$&$n$&$\binom{n-1}2$&$4$\\
$r=2,p=2,n\text{ odd}$&$W_2$&$\binom{n-1}2$&$n-2$&$4$\\\bottomrule
\end{tabular}
\end{center}
In every row $\dim C_r=N-\dim RC_r$.
The reflexive duality in all five rows follows from the common scalar
classification; it is not a retained computational assertion.
The final row, including both distances and both dimensions, is proved
by \Uref{ULCD.Johnson} and \Uref{ULCD.SimplexMin}.
In the third row \Uref{ULCD.TernaryKneser} proves the space identity and
dimension: $\dim W_2=\binom n2-1$ and $\mathbf j\notin W_2$.
Thus the two printed dimensions $\binom n2+1$ and
$N-\binom n2-1$ need the indicated one-unit correction.
The final row similarly corrects the printed $\binom n2$ to
$\binom{n-1}2$.
The remaining space identifications and distances in the first four
rows are retained, exactly as quoted by the source from its earlier
papers; the common projector proof alone is not a distance proof.
This distinction includes the distance $8$ in the ternary row.
The first row's upper bound $8$ also follows directly from a nonzero
alternating transversal word on six points, using $W_\Pi\subseteq W_2^\perp$.

For $k=2,n\geq5,p=2$, source Result~3 has two additional cases.
If $n\equiv3\pmod4$, the identification
$RC_0=W_1+\langle\mathbf j\rangle$ has parameters
$[\binom n2,n,n-1]_2$, with dual
$[\binom n2,\binom n2-n,4]_2$.
Indeed $RC_0=\im(J+A_{J(n,2)})$; the two summands are annihilating
projectors onto $\langle\mathbf j\rangle$ and $W_1$, respectively,
since $\binom n2$ is odd and the Johnson degree is even.
Their images therefore add directly. The two distances just stated
remain the source's quoted earlier-code results.
For either odd residue of $n$, the Johnson code is
$C_1=W_1$, with parameters $[\binom n2,n-1,n-1]_2$;
its reflexive dual has parameters
$[\binom n2,\binom n2-n+1,3]_2$.
The missing $+1$ in the printed dual dimension is corrected by the
simplex ranks. These Johnson parameters and distances are fully proved
in the common paper.
If $n\equiv1\pmod4$, the quoted $C_0=E(W_1)^\perp$ is LCD.
Here $E(W_1)=W_1$: the coefficient kernel is the odd-weight constant
vector, so every coefficient class has an even-weight representative.
The common complement theorem, applied to the complementary Johnson
graph whose code is $W_1$, gives $C_0=W_1^\perp$ and proves LCD.

Source Result~5 is a separately quoted incidence theorem: for
$n\geq2k+1$, $W_{k-1}$ has minimum weight $n-k+1$ over every prime
field; for $n>2k+1$ its minimum words are scalar multiples of the stars
$w_\Lambda$; its dimension is at least $\binom{n-1}{k-1}$.
The binary equality, minimum weight and minimum-word statement are
proved by \Uref{ULCD.Simplex} and \Uref{ULCD.SimplexMin}.
The odd-characteristic clauses remain retained source reference [11]
inputs. Likewise the minimum $\binom{n-1}2$ for the binary $W_1$ on
triples is retained from the quoted Result~4; the exact enumerator above
gives an alternative reduction to an integer inequality but is not
silently substituted for that inequality.

\subsection{All four Johnson parity cases and decoding}
Put $E=u_{k-1}d_k$ on binary $k$-subsets, in the notation of the common
simplex section. The identity $A=E+(k\bmod2)I$, with $E^2=nE$, proves
all four cases of source Result~6. For odd $n$, $E$ is the projector
onto $W_{k-1}$: $C=W_{k-1}^\perp,RC=W_{k-1}$ for odd $k$, and the
two are interchanged for even $k$. Both hulls vanish.
For even $n$, $\im E\subseteq\operatorname{Hull}(W_{k-1})$;
the adjacency and reflexive codes are $(\im E,V_k)$ for even $k$ and
$(V_k,\im E)$ for odd $k$. These are consequences of
\Uref{ULCD.Johnson}, without the source's earlier-code citation as an
extra premise.

In source Example~1, $n>7$, $n\equiv0\pmod4$, let
$C=RC_1=W_1$. For a received binary vector $v$ with support $S$,
\[
 v=\sum_{x\in S}s_x^1+\sum_{x\in S}r_x^1=v_C+v_\perp.
\]
The preceding quoted parameters give minimum distance
$\binom{n-1}2$. A coset leader $c$ of $C-v_\perp$ yields the decoded
word $v+c$, uniquely if fewer than $\tfrac12\binom{n-1}2$ errors
occurred. This is exactly the nearest-word argument of the projection
section. The source's suggestion that a faster explicit decoder would
be desirable is recorded as a question; no new decoding-complexity claim
is inferred from this decomposition.
